\section{Score Coordinates and Answer Discrimination}
\label{app:response-diagnostics}

We retrospectively analyze the stored outputs from the 400-patient Qwen
causal-ownership matrix cohort at $\alpha=+0.25$, distinct from the earlier
400-patient Effusion specificity supplement. The logit margin is the yes-token
logit minus the no-token logit; each intervention effect is its change
relative to the same patient's unperturbed output. In each score coordinate,
the ownership margin subtracts the largest mean effect among the five other
clinical directions from the matched direction's mean effect.
Table~\ref{tab:response-diagnostics} compares these margins with answer
discrimination. Here, AUROC ranks patients by the model's answer logit
margin against the corresponding disease labels, rather than by a fitted
probe score. The probability-ownership values here belong to the matrix
cohort: its Effusion margin is $-0.0640$, whereas the earlier specificity
supplement reports $-0.0573$ on different patients.

\begin{table}[H]
\centering
\caption{Score-coordinate ownership and answer discrimination.}
\label{tab:response-diagnostics}
\resizebox{\linewidth}{!}{%
\begin{tabular}{lrrrr}
\toprule
Question & \shortstack{Probability\\ownership} & \shortstack{Logit-margin\\ownership} & \shortstack{Baseline\\AUROC} & \shortstack{Matched-direction\\AUROC} \\
\midrule
Effusion     & $-0.0640$ & $-0.4853$ & 0.6178 & 0.5921 \\
Atelectasis  & $-0.2295$ & $-1.0703$ & 0.6235 & 0.5479 \\
Pneumothorax & $-0.1086$ & $-0.8944$ & 0.6666 & 0.6324 \\
Cardiomegaly & $-0.1972$ & $-0.8559$ & 0.6217 & 0.5152 \\
Mass         & $-0.1528$ & $-0.7134$ & 0.6636 & 0.6354 \\
Nodule       & $-0.0831$ & $-0.4134$ & 0.6805 & 0.6355 \\
\bottomrule
\end{tabular}%
}
\end{table}

All six ownership margins remain negative before the sigmoid mapping to
probability. Answer-score AUROC is lower under each matched direction in
the point estimates, separating a larger affirmative response from better
patient discrimination. These are descriptive comparisons, with the
registered ownership decisions retaining their original inferential status.

\subsection{Label-conditioned shifts and matrix geometry}

For the Effusion direction on the Effusion question, the mean logit-margin
shift is 2.1336 among 29 label-positive patients and 2.2712 among 371
label-negative patients. The positive-minus-negative difference is
$-0.1376$, with a pointwise 95\% patient-bootstrap percentile interval of
$[-0.4443,0.1659]$ from 2,000 draws. Both groups show an upward shift,
while their mean-shift difference remains unresolved. The Brier score,
the mean squared difference between the yes probability and binary label,
rises from 0.0681 to 0.1035.

The six-by-six matrix of mean logit-margin changes has clinical directions
as rows and questions as columns. Its first singular component accounts
for 90.51\% of the sum of squared matrix entries. Subtracting row and column
means and restoring the grand mean leaves an interaction component with
7.91\% of that same raw matrix energy. These are distinct decompositions
with a common denominator, not complementary percentages. They summarize
response geometry rather than identifying a shared causal channel. The
answer-encoding experiment below tests another feature of that response
structure: its behavior when the finding-to-letter mapping is exchanged.

\FloatBarrier

\section{Answer-Encoding Sensitivity}
\label{app:encoding}

\subsection{Paired design and independent capability calibration}

The crossover selects the first 200 patients in the accepted causal-ownership
matrix cohort's registered order for interventions and the remaining 200 for
capability calibration. Each patient contributes one image. These two
sets are patient-disjoint from each other but both reuse the original
ownership cohort. The six clinical questions cross three prompt packages,
each question--prompt pair having its own unsteered baseline and
27 directions at each of $\alpha=\pm0.25$:
six clinical directions, 20 seed-0 random directions, and a coordinate-
permutation sham. The 198,000 intervention outcomes and 2,400 clean
calibration outcomes took 6,402 seconds on one A100.

Figure~\ref{fig:encoding-prompts} displays the exact Mass prompts. The other
five questions use the corresponding finding phrases from the original
clinical family. Coded prompts exchange the answer mapping while retaining
their wording; comparison with standard yes/no changes the question wording
and response instructions as well as the answer symbols.

\begin{figure}[H]
\centering
\begin{minipage}{0.95\linewidth}
\small
\fbox{\parbox{0.95\linewidth}{Is there a lung mass in this chest radiograph? Answer yes or no.}}
\par\smallskip\centerline{(a) Standard yes/no}\medskip
\fbox{\parbox{0.95\linewidth}{Does this chest radiograph show a lung mass? Answer A if the finding is present and B if it is absent. Reply with A or B only.}}
\par\smallskip\centerline{(b) A-present}\medskip
\fbox{\parbox{0.95\linewidth}{Does this chest radiograph show a lung mass? Answer B if the finding is present and A if it is absent. Reply with A or B only.}}
\par\smallskip\centerline{(c) B-present}
\end{minipage}
\caption{\textbf{The crossover distinguishes answer mapping within coded prompts from the full prompt-package comparison.} The two coded conditions exchange presence and absence assignments for A and B. The standard condition retains the original question.}
\label{fig:encoding-prompts}
\end{figure}

At the first answer position, the raw coded margin is the maximum A-token
logit minus the maximum B-token logit among valid single-token leading-space
variants. A-present uses that margin as the finding-present score;
B-present negates it. Capability calibration uses these clinically oriented
scores against image labels. Both mappings must have a one-sided 95\%
AUROC lower bound above 0.5 and at least ten patients in each label class.
The lower bounds use 2,000 patient-bootstrap draws, with seed 20260908.
Table~\ref{tab:encoding-calibration} reports the resulting eligibility.

\begin{table}[H]
\centering
\caption{Independent answer-encoding capability calibration.}
\label{tab:encoding-calibration}
\resizebox{\linewidth}{!}{%
\begin{tabular}{lrrrrrl}
\toprule
Question & Positive & \shortstack{A-present\\AUROC} & \shortstack{B-present\\AUROC} & \shortstack{A-present\\lower bound} & \shortstack{B-present\\lower bound} & Eligibility \\
\midrule
Effusion     & 11 & 0.5125 & 0.5111 & 0.3872 & 0.3967 & Below baseline criterion \\
Atelectasis  & 15 & 0.6586 & 0.6441 & 0.5448 & 0.5372 & Eligible \\
Pneumothorax & 13 & 0.6508 & 0.6524 & 0.5088 & 0.5113 & Eligible \\
Cardiomegaly &  8 & 0.5661 & 0.5811 & 0.4677 & 0.4757 & Insufficient labels \\
Mass         & 15 & 0.6533 & 0.6542 & 0.5198 & 0.5231 & Eligible \\
Nodule       & 21 & 0.6347 & 0.6327 & 0.5210 & 0.5196 & Eligible \\
\bottomrule
\end{tabular}%
}
\end{table}

Atelectasis, Pneumothorax, Mass, and Nodule enter the primary endpoint.
Effusion and Cardiomegaly remain descriptive. Clean finding-present A/B
rankings are highly correlated: independent calibration Spearman
correlations range from 0.960770 to 0.988116, while the intervention-cohort
range is 0.9594--0.9804. Calibration mean raw A-minus-B margins have
opposite signs across mappings for every question; their absolute
magnitudes differ by 0.0050--0.1200. Thus similar patient rankings coexist
with baseline score asymmetry.

\subsection{Mapping-even and mapping-odd response energy}

For patient $i$, direction $d$, and question $q$, let $x^A_{i,d,q}$ and
$x^B_{i,d,q}$ denote the changes in raw A-minus-B margin from each prompt's
own baseline under A-present and B-present, respectively, at
$\alpha=+0.25$. Define the mapping-even component $e_{i,d,q}$ and
mapping-odd component $o_{i,d,q}$ by
\begin{equation}
e_{i,d,q}=\tfrac12(x^A_{i,d,q}+x^B_{i,d,q}),\qquad
o_{i,d,q}=\tfrac12(x^A_{i,d,q}-x^B_{i,d,q}).
\label{eq:encoding-components}
\end{equation}
The even component retains raw letter preference under the exchange; the
odd component reverses raw sign as the clinical assignment reverses.
These are coordinate decompositions, not mechanism labels.

\begin{samepage}
For a sequence of $n=200$ patient effects $z=(z_1,\ldots,z_n)$, let
$U(z)$ estimate its squared population mean. Let $\mathcal Q$ be the four
eligible questions and $\mathcal D$ the six clinical directions. The
registered energy contrast $T$ is
\begin{equation}
U(z)=\frac{(\sum_i z_i)^2-\sum_i z_i^2}{n(n-1)},\qquad
T=\frac{1}{6}\sum_{d\in\mathcal D}\sum_{q\in\mathcal Q}
\bigl[U(e_{\cdot,d,q})-U(o_{\cdot,d,q})\bigr].
\label{eq:encoding-energy}
\end{equation}
The cross-patient products remove the finite-sample variance contribution
to a squared sample mean.
\end{samepage}
The endpoint sums questions and averages clinical
directions, in squared-logit units. Joint patient resampling preserves
pairing across mappings, directions, and questions; 2,000 draws use seed
20260907. Each random direction uses the same eligible-question sum.
For sham, the question-specific permuted normal supplies each question's
effect. The registered even-beyond-controls condition requires a positive
95\% lower endpoint for $T$, a point estimate above every random contrast,
and a point estimate above the absolute sham contrast.

Table~\ref{tab:encoding-energy} compares the primary estimate with its
controls and the all-question descriptive endpoint. Negative values favor
mapping-odd squared mean response in the aggregate.

\begin{table}[H]
\centering
\caption{Mapping-even versus mapping-odd response energy.}
\label{tab:encoding-energy}
\resizebox{\linewidth}{!}{%
\begin{tabular}{llrr}
\toprule
Comparison & Questions & Estimate & 95\% interval \\
\midrule
Clinical mean & Four eligible & $-2.1786$ & $[-2.4061,-1.9810]$ \\
Random-direction range & Four eligible & $[-5.4773,-0.1769]$ & --- \\
Sham & Four eligible & $-2.3848$ & $[-2.5371,-2.2460]$ \\
Clinical mean & All six & $-2.9329$ & $[-3.2374,-2.6636]$ \\
\bottomrule
\end{tabular}%
}
\end{table}

The primary interval is entirely negative, and all 20 random-direction
point estimates are negative as well. The registered even-beyond-controls
condition is therefore unmet. Mapping-odd dominance is a property shared
with generic perturbations in this experiment, so the energy contrast does
not by itself establish clinical specificity or encoding invariance. The
random row gives the observed range over directions, not an uncertainty
interval; the other intervals are patient-bootstrap percentile summaries.

\subsection{Prompt-conditioned clinical competition}

Table~\ref{tab:encoding-mass} reports Mass in finding-present logit
coordinates, using yes-minus-no for standard prompts, A-minus-B for
A-present, and B-minus-A for B-present. Each clinical margin subtracts the
maximum alternative among five directions within that same prompt; the
maximum is recomputed in each patient-bootstrap draw. The matched effects
under A-present and B-present are 1.6125 and 1.7469, exceeding their random
maxima of 1.1650 and 1.2581 and absolute sham effects of 0.7550 and 1.1475.
Their positive clinical margins contrast with the standard value of
$-0.7181$ on the same 200 patients. That standard-prompt value is recomputed
on the intervention subset; the full 400-patient diagnostic value is
$-0.7134$ (Table~\ref{tab:response-diagnostics}).

This is a descriptive discovery selected from six questions, with
pointwise intervals and 20 random directions. Standard-to-coded contrasts
concern the whole prompt package because wording and response instructions
also change. The result motivates independent, prospectively locked
confirmation of prompt-conditioned specificity; it preserves the original
400-patient simultaneous ownership verdict at its original yes/no prompts.

\FloatBarrier
